\documentclass[10pt,a4paper]{amsart}
\usepackage[margin=19mm]{geometry}
\usepackage{amsmath,amssymb,mathtools}
\usepackage[scr=rsfs,cal=euler]{mathalfa}
\usepackage{xfrac}
\usepackage[unicode,hidelinks,bookmarks=false]{hyperref}
\newcommand{\ie}{i.e.\ }
\newcommand{\cf}{cf.\ }
\newcommand{\resp}{respectively\ }
\newcommand{\Subsection}{\subsection}
\providecommand{\nobreakdash}{\nobreak\hbox{-}}
\setlength{\emergencystretch}{2em}
\multlinegap0pt
 \newcommand{\Z}{{\mathbb Z}}
\usepackage{amssymb}
\usepackage{dsfont}
\usepackage{tikz}
\usepackage{bm}
\usepackage{mathtools}
\newtheorem{proposition}{Proposition}
\title{The functor between two categories of $\Z-$graded manifolds}
\author{Martha Valentina Guarin Escudero, Alexei Kotov}
\date{Source: arXiv:2602.02420v2 (CC BY 4.0)}
\begin{document}
\maketitle

\noindent\emph{Source and licence.} The functor between two categories of $\Z-$graded manifolds. Version arXiv:2602.02420v2 (CC BY 4.0), distributed by arXiv under the Creative Commons Attribution 4.0 International licence, http://creativecommons.org/licenses/by/4.0/, as recorded in the arXiv licence metadata for this exact version and shown on the abstract page https://arxiv.org/abs/2602.02420v2. Full licence notice: \texttt{licenses/arxiv-2602.02420-CC-BY-4.0.txt}.\par\smallskip

\noindent\textit{Editorial scope.} Counted unit: the article's proposition prop:graded-completion (source0.tex lines 1193--1195). The article's complete original proof of it is delivered with it (source0.tex lines 1197--1210, one \texttt{proof} environment of 14 lines). No mathematics is written, repaired or re-derived: the statement and the proof are the article's own lines, in the article's own notation, and no retained line is substituted or re-flowed.  No further source-local context is needed: every cross-reference of every retained line resolves inside the two delivered spans.  Nothing was omitted from any retained span, so the package carries no omission record.  No retained line contains a citation, so the wrapper carries no bibliography and no bibliography entry.  No retained line contains a figure environment, an image-inclusion command or drawing code, so no image is delivered and the package declares no asset dependency.  Editorial formatting only, in addition: an AMS \texttt{amsart} preamble that restates the article's own declarations and macros used by the retained lines, namely the environments \texttt{proposition} on separate counters, together with the standard packages those lines need.  The retained environments print in this document as Proposition 1 in order of appearance, whereas the article prints the same items with the numbering of its own full document.  The complete native source of the article as distributed in the arXiv e-print for this exact version is bundled unchanged as \texttt{sources/193851/source0.tex}.
\begin{proposition}[Graded completions]\label{prop:graded-completion}\normalfont
A $\mathbb{Z}$-graded filtered algebra $A$ admits a \emph{graded completion} with respect to the graded filtration, given componentwise by $\widehat{A} = \bigoplus_{\gamma\in\Gamma} \widehat{A}_\gamma$, where $\widehat{A_\gamma}=\varprojlim_i A_\gamma/(F^i A_\gamma)$. Equivalent filtrations (on each $A_\gamma$ independently) yield isomorphic graded completions.
\end{proposition}

\noindent\begin{proof}
In a $\Gamma$-graded algebra $A$, the homogeneous multiplications $A_{\gamma_1}\times A_{\gamma_2}\to A_{\gamma_1+\gamma_2}$ are continuous in the filtration topology and thus extend to $\widehat{A}_{\gamma_1}\times\widehat{A}_{\gamma_2}\to\widehat{A}_{\gamma_1+\gamma_2}$.

\noindent
Indeed, let $(a_n)_{n\in\mathbb{N}}$, $(a'_n)_{n\in\mathbb{N}}$ be Cauchy sequences in $A_{\gamma_1}$, $A_{\gamma_2}$. For each $k$, there exist $n(k)$, $n'(k)$ such that
\[
a_n\equiv a_{n(k)}\pmod{F_kA_{\gamma_1}}\ \forall n\ge n(k), \quad a'_n\equiv a'_{n'(k)}\pmod{F_kA_{\gamma_2}}\ \forall n\ge n'(k).
\]
Their product $(a_na'_n)_n$ is Cauchy:
\[
a_na'_n\equiv a_{n(k)}a'_{n'(k)}\pmod{F_kA_{\gamma_1+\gamma_2}}\ \forall n\ge\max(n(k),n'(k)).
\]
Moreover, this respects Cauchy equivalence classes, yielding a well-defined $\Gamma$-graded completion of $A$.  $\square$
\end{proof}

\end{document}
